% !TeX program = lualatex
% !TeX encoding = utf8
% !TeX spellcheck = uk_UA
% !TeX root = ../EMConspect.tex
% УВАГА: для \enquote{...} потрібен пакет csquotes (\usepackage[ukrainian=quotes]{csquotes})
% Якщо в преамбулі ще не визначено \Mfield для вектора намагніченості, безпечно визначаємо тут:
\providecommand{\Mfield}{\vect{M}}

%=========================================================
\chapter{Магнітне поле в речовині}\label{\currfilebase}
\graphicspath{{\currfilebase/Pictures}}
%=========================================================
\epigraph{Речовина в магнітному полі, як і в електричному, --- не пасивне тло. Різниця лише в тому, \emph{чим} вона відповідає: діелектрик
	зміщує заряди, а магнетик --- перебудовує вже наявні в ньому мікроскопічні струми.
}{Переосмислення ідей сучасної фізики}

У попередньому розділі ми вивчали магнітне поле джерел --- провідників зі струмом --- у \emphx{вакуумі}. Проте більшість тіл, з якими має справу
техніка (і які трапляються в природі), --- це не вакуум: провід у котушці ізольований лаком, осердя трансформатора зроблене із заліза, компас
працює в повітрі, а зразок для дослідження часто занурюють у рідину. У всіх цих випадках магнітне поле існує \emphx{в присутності речовини}, і
речовина на це поле відповідає --- \emphx{намагнічується}. Тут ми побачимо, що ця відповідь --- за формою --- напрочуд схожа на поляризацію
діелектриків (гл.~\ref{FileldInDielectrics}), хоча її фізична природа --- зовсім інша: замість зміщення зв'язаних \emphx{зарядів} тут працює
перебудова вже наявних \emphx{мікрострумів}.

%% --------------------------------------------------------
\section{Магнетики: вільні й молекулярні струми}
%% --------------------------------------------------------

Точнісінько так само, як усі заряди речовини ми умовно розділяли на вільні й зв'язані (гл.~\ref{FileldInDielectrics}), усі \emphx{струми} в
речовині зручно розділити на два типи.

\begin{enumerate}
	\item \emphx{Вільні (макроскопічні) струми} $\vect j$ --- це струми, якими ми керуємо ззовні: струм провідності в дротах котушки, струм у
	      плазмі тощо. Саме такі струми ми задаємо чи вимірюємо в лабораторії.
	\item \emphx{Молекулярні (зв'язані) струми} $\vect j'$ --- це мікроскопічні струми, що циркулюють \emphx{у межах кожного окремого атома чи
	      молекули} й ніколи не виходять за ці межі: рух електронів навколо ядра орбітою, а також їхній власний (спіновий) магнетизм. Кожен такий
	      мікрострум, за означенням~\eqref{eq:MagMomentLin}, еквівалентний крихітному магнітному диполю --- \emphx{магнітному моменту атома}.
\end{enumerate}

\begin{Attention}
	\emphx{Магнетиком} називають будь-яку речовину: адже, на відміну від діелектриків (де зв'язаних зарядів може взагалі не бути в
	провідниках), молекулярні струми --- а отже, і магнітні моменти атомів --- є \emphx{в кожному без винятку} атомі: рух електронів навколо ядра
	не припиняється ніколи. Різні речовини відрізняються не наявністю чи відсутністю цих мікрострумів, а тим, \emphx{як} вони поводяться в
	зовнішньому полі і \emphx{чи компенсують} вони одне одного за його відсутності.
\end{Attention}

Так само, як і для електричного поля, справжнє (мікроскопічне) магнітне поле всередині речовини надзвичайно неоднорідне: поблизу ядра воно
величезне, а між атомами --- слабке, і швидко змінюється в часі разом із рухом електронів. Макроскопічні прилади фіксують лише поле, усереднене за
фізично малим об'ємом --- саме таке усереднене поле ми й позначатимемо надалі просто $\Bfield$, коли йтиметься про поле в речовині.

%% --------------------------------------------------------
\section{Магнітні моменти атомів. Класифікація магнетиків}
%% --------------------------------------------------------

\subsection*{Орбітальний і спіновий магнетизм електрона}

Електрон, що рухається навколо ядра, --- це, по суті, крихітний виток зі струмом, а отже, володіє \emphx{орбітальним} магнітним моментом,
пов'язаним з орбітальним моментом імпульсу $\vect L_{\text{орб}}$ через гіромагнітне відношення~\eqref{eq:GyromagneticClassical} (гл.~\ref{MagFieldVacuum}):
\begin{equation}\label{eq:OrbitalMoment}
	\vect p_{\text{орб}} = -\frac{e}{2m_ec}\, \vect L_{\text{орб}},
\end{equation}
де знак \enquote{$-$} відображає від'ємний заряд електрона: орбітальний магнітний момент \emphx{антипаралельний} орбітальному моменту імпульсу.

Крім орбітального руху, електрон має власний (спіновий) момент імпульсу $\vect S$ і пов'язаний з ним \emphx{спіновий} магнітний момент, --- той
самий \enquote{аномальний} результат, на який ми вже звертали увагу (Увага після формули~\eqref{eq:GyromagneticClassical}):
\begin{equation}\label{eq:SpinMoment}
	\vect p_{\text{спін}} = -\frac{e}{m_ec}\, \vect S.
\end{equation}
\emphx{Повний магнітний момент атома} --- це векторна сума орбітальних і спінових моментів усіх його електронів (внесок ядра в
$10^3$--$10^4$ разів менший через велику масу протона і тому практично завжди нехтовний):
\begin{equation}\label{eq:AtomTotalMoment}
	\vect p_{\text{атом}} = \sum_i \left(\vect p_{\text{орб},i} + \vect p_{\text{спін},i}\right).
\end{equation}

\subsection*{Класифікація магнетиків}

Ключове питання --- чи дорівнює ця сума нулю за \emphx{відсутності} зовнішнього поля. Відповідь залежить від електронної будови атома чи
молекули, і саме за цією ознакою всі речовини поділяють на класи, цілком аналогічні до полярних і неполярних діелектриків.

\begin{itemize}
	\item \emphx{Діамагнетики} --- речовини, атоми (молекули) яких мають \emphx{повністю заповнені} електронні оболонки. Орбітальні й спінові
	      моменти окремих електронів у такій оболонці попарно компенсують одне одного (кожному електрону з проекцією моменту \enquote{вгору} відповідає
	      \enquote{напарник} з проекцією \enquote{вниз}), тому власний магнітний момент атома в цілому дорівнює нулю: $\vect p_{\text{атом}} = 0$. Це
	      \emphx{точний аналог неполярної молекули} ($p_0 = 0$) з гл.~\ref{FileldInDielectrics}.
	\item \emphx{Парамагнетики} --- речовини, атоми (молекули) яких мають \emphx{незаповнені} електронні оболонки (непарну кількість
	      електронів, або оболонку з нескомпенсованим спіном чи орбітальним моментом). Такий атом має власний, ненульовий магнітний момент
	      $\vect p_{\text{атом}} \neq 0$ навіть без поля, --- \emphx{точний аналог полярної молекули} ($p_0 \neq 0$).
	\item \emphx{Феромагнетики (і споріднені їм феримагнетики, антиферомагнетики)} --- особливий, значно складніший клас речовин, у яких
	      сусідні атомні моменти не лише існують, а й \emphx{самочинно} (без зовнішнього поля) вишиковуються паралельно одне одному внаслідок
	      квантової \emphx{обмінної взаємодії}. Це --- прямий магнітний аналог сегнетоелектриків (гл.~\ref{FileldInDielectrics}, п.~\ref{sec:FE}):
	      і там, і там причина --- не зовнішнє поле, а внутрішня, кооперативна взаємодія сусідніх диполів. Через істотну відмінність механізму
	      (доменна структура, гістерезис, точка Кюрі) феромагнетики заслуговують на окремий розділ і тут детально не розглядаються.
\end{itemize}

\begin{table}[h!]\small\centering
	\caption{Класифікація магнетиків за поведінкою атомного магнітного моменту\label{tab:MagClassification}}
	\begin{tabular}{lll}
		\toprule
		Клас             & Власний момент атома      & Приклади                                                              \\
		\midrule
		Діамагнетики     & $\vect p_{\text{атом}}=0$ & інертні гази, $\ce{N2}$, $\ce{H2}$, вода, більшість органіки, $\ce{Cu}$, $\ce{Au}$, $\ce{Bi}$ \\
		Парамагнетики    & $\vect p_{\text{атом}}\neq0$ & $\ce{O2}$, лужні й лужноземельні метали, $\ce{Al}$, $\ce{Pt}$, солі перехідних металів \\
		Феромагнетики    & $\vect p_{\text{атом}}\neq0$, впорядковані & $\ce{Fe}$, $\ce{Co}$, $\ce{Ni}$, деякі сплави й рідкісноземельні метали \\
		\bottomrule
	\end{tabular}
\end{table}

\begin{Attention}
	Наявність власного моменту --- \emphx{необхідна}, але не єдина умова: у парамагнетика цей момент, за відсутності поля, орієнтований цілком
	хаотично (тепловий рух), і в середньому по речовині компенсується --- так само, як дипольні моменти полярних молекул у неполяризованому
	діелектрику. \emphx{Намагніченість} з'являється лише тоді, коли зовнішнє поле частково впорядковує ці моменти. У феромагнетика ж
	впорядкування виникає \emphx{самочинно}, ще до накладання поля, --- у цьому й полягає його принципова відмінність.
\end{Attention}

%% --------------------------------------------------------
\section{Природа діамагнетизму}\label{sec:DiaNature}
%% --------------------------------------------------------

Почнемо з діамагнетизму --- явища, яке, на відміну від парамагнетизму й феромагнетизму, властиве \emphx{абсолютно всім} речовинам без винятку
(включно з парамагнетиками й феромагнетиками!), оскільки його джерело --- сам факт руху електронів по орбітах, а не наявність чи відсутність
власного атомного моменту.

\subsection*{Якісна картина: індукований момент проти поля}

Уявімо електрон, що обертається орбітою навколо ядра. Увімкнемо зовнішнє поле $\Bfield$, перпендикулярне площині орбіти. За законом
електромагнітної індукції (з яким ми детально познайомимось пізніше, але результат якого тут можна передбачити якісно), \emphx{зростання} поля
$\Bfield$ від нуля до кінцевого значення саме по собі створює вихрове електричне поле, яке гальмує чи прискорює електрон --- залежно від напрямку
обертання --- так, щоб, за правилом Ленца, \emphx{перешкодити} зростанню потоку, тобто \emphx{послабити} прикладене поле. У результаті орбітальний
рух електрона отримує невелику додаткову складову, яка відповідає магнітному моменту, спрямованому \emphx{проти} поля $\Bfield$, --- \emphx{незалежно
від того, у який бік електрон обертався початково}.

\begin{Attention}
	Це --- принципова відмінність від електронної поляризації діелектриків (гл.~\ref{FileldInDielectrics}), де наведений дипольний момент
	з'являється через просте \emphx{зміщення} електронної хмари вздовж поля $\Efield$ і тому напрямлений \emphx{за} полем. Тут же індукований
	магнітний момент виникає через зміну \emphx{швидкості обертання} електрона і завжди напрямлений \emphx{проти} поля --- це universаль\-ний
	прояв закону збереження \enquote{енергії заважати змінам}, тобто прояв правила Ленца на мікроскопічному, атомному рівні.
\end{Attention}

\subsection*{Ларморова прецесія орбіти}

Дамо цій картині кількісний вигляд. Орбітальний магнітний момент електрона пов'язаний з його орбітальним моментом імпульсу
формулою~\eqref{eq:OrbitalMoment}: $\vect p_{\text{орб}} = \gamma\vect L_{\text{орб}}$, де $\gamma = -e/(2m_ec)$. У зовнішньому полі $\Bfield$ на
цей магнітний момент діє момент сили~\eqref{eq:TorqueDipoleAnalogy}, який, за основним рівнянням динаміки обертального руху, змінює момент
імпульсу:
\begin{equation*}
	\frac{d\vect L_{\text{орб}}}{dt} = \vect p_{\text{орб}}\times\Bfield = \gamma\left(\vect L_{\text{орб}}\times\Bfield\right) = -\gamma\left(\Bfield\times\vect L_{\text{орб}}\right).
\end{equation*}
Це --- рівняння виду $d\vect L_{\text{орб}}/dt = \vect\omega_L\times\vect L_{\text{орб}}$ із $\vect\omega_L \equiv -\gamma\Bfield$, а
таке рівняння описує не що інше, як \emphx{рівномірну прецесію} вектора $\vect L_{\text{орб}}$ навколо напрямку поля $\Bfield$ з кутовою
швидкістю
\begin{equation}\label{eq:LarmorFreq}
	\tcbhighmath{
		\omega_L = |\gamma| B = \frac{eB}{2m_ec}.
	}
\end{equation}
Величину $\omega_L$ називають \emphx{ларморовою частотою прецесії}, а саме явище --- \emphx{ларморовою прецесією}.

\begin{Attention}
	Фізичний зміст ларморової прецесії такий: у зовнішньому полі $\Bfield$ уся площина електронної орбіти (а разом з нею --- і сам вектор
	орбітального моменту імпульсу) починає повільно обертатись навколо напрямку $\Bfield$ з частотою $\omega_L$, \emphx{додатково} до швидкого
	обертання електрона по самій орбіті. Важливо: частота $\omega_L$ \eqref{eq:LarmorFreq} \emphx{не залежить} ні від радіуса орбіти, ні від
	швидкості електрона на ній, ні від напрямку його початкового обертання, --- а залежить лише від поля $\Bfield$ і від відношення заряду до
	маси електрона $e/m_e$. Це --- \emphx{теорема Лармора}: у слабкому магнітному полі рух будь-якої системи зарядів з однаковим відношенням
	$e/m$ можна отримати з руху \emphx{без поля}, додавши до нього прецесію із частотою $\omega_L$ навколо напрямку $\Bfield$.
\end{Attention}

\subsection*{Індукований діамагнітний момент}

Прецесія орбіти з кутовою швидкістю $\omega_L$ --- це, у свою чергу, додатковий круговий рух заряду $-e$ навколо осі, паралельної $\Bfield$, тобто
додатковий \emphx{круговий струм}. Нехай $\rho$ --- відстань від електрона до осі прецесії (осі, що проходить через ядро паралельно $\Bfield$).
Тоді цей додатковий струм дорівнює заряду, помноженому на частоту обертання, $I = -e\,\omega_L/(2\pi)$, а відповідний йому додатковий магнітний
момент --- за формулою~\eqref{eq:MagMomentLoop} для колового витка, $p = IS/c$ із площею $S=\pi\rho^2$:
\begin{equation}\label{eq:InducedDiaMoment}
	\Delta p = \frac{I\pi\rho^2}{c} = -\frac{e\,\omega_L\,\rho^2}{2c} = -\frac{e^2 B \rho^2}{4m_ec^2},
\end{equation}
де ми підставили $\omega_L$ із~\eqref{eq:LarmorFreq}. Знак \enquote{$-$} тут не випадковий: він означає, що $\Delta p$ спрямований \emphx{проти}
поля $\Bfield$ незалежно від того, як електрон рухався початково, --- рівно те, що ми передбачили якісно.

Усереднюючи $\rho^2 = x^2+y^2$ (де вісь прецесії --- це напрямок $\Bfield$) за сферично-симетричним розподілом електронної хмари в атомі, для
якого $\langle x^2\rangle = \langle y^2\rangle = \langle z^2\rangle = \langle r^2\rangle/3$, маємо $\langle\rho^2\rangle = \langle
r^2\rangle - \langle z^2\rangle= \frac23\langle r^2\rangle$. Підсумовуючи внесок усіх $Z$ електронів атома, отримуємо \emphx{індукований
діамагнітний момент атома}:
\begin{equation}\label{eq:AtomDiaMoment}
	\tcbhighmath{
		\vect p_{\text{діа}} = -\frac{Z e^2\langle r^2\rangle}{6m_ec^2}\, \Bfield,
	}
\end{equation}
де $\langle r^2\rangle$ --- середній квадрат відстані електронів від ядра (характерний атомний масштаб, $\langle r^2\rangle^{1/2}\sim 1$~\AA).

\subsection*{Формула Ланжевена. Сприйнятливість діамагнетика}

Якщо в одиниці об'єму міститься $n$ атомів, то намагніченість (див. означення нижче, формула~\eqref{eq:MagnetizationDef}) дорівнює $\Mfield = n\,
\vect p_{\text{діа}}$, а оскільки в такому слабкому й лінійному відгуку різниця між $\Bfield$ і $\Hfield$ у самому магнетику неістотна (поправка
вищого порядку малості за $\chi$), можна записати $\Mfield \approx \chi\Bfield$, звідки
\begin{equation}\label{eq:LangevinDia}
	\tcbhighmath{
		\chi_{\text{діа}} = -\frac{nZe^2\langle r^2\rangle}{6m_ec^2}.
	}
\end{equation}
Цей результат називають \emphx{формулою Ланжевена для діамагнетизму}. Зверніть увагу на її головні риси, які повністю узгоджуються з
експериментом:

\begin{itemize}
	\item $\chi_{\text{діа}} < 0$ завжди --- діамагнітний відгук завжди \emphx{послаблює} зовнішнє поле всередині речовини (магнітний аналог
	      правила Ленца в статиці);
	\item $\chi_{\text{діа}}$ \emphx{не залежить від температури} --- на відміну від парамагнітної сприйнятливості (див. далі), оскільки
	      розмір атома $\langle r^2\rangle$ практично не змінюється з нагріванням;
	\item за порядком величини $\chi_{\text{діа}}\sim 10^{-5}$--$10^{-6}$ --- діамагнетизм є \emphx{дуже слабким} ефектом;
	\item діамагнітний внесок $\chi_{\text{діа}}$ присутній \emphx{в кожному атомі}, включно з парамагнітними й феромагнітними, --- просто в
	      них він на порядки менший за власний (орієнтаційний чи обмінний) внесок і тому непомітний на його тлі.
\end{itemize}

\begin{table}[h!]\small\centering
	\caption{Магнітна сприйнятливість деяких діамагнетиків (кімнатна температура)\label{tab:ChiValues}}
	\begin{tblr}{
		colspec={X[l, m]X[c,m]},
		row{1} = {c, m, fg=white, bg=cyan, font=\bfseries},
		hlines, vlines
	}
		Речовина & $\chi\ (\times10^{-6})$ \\
		Вода ($\ce{H2O}$) & $-9$ \\
		Мідь ($\ce{Cu}$) & $-10$ \\
		Золото ($\ce{Au}$) & $-28$ \\
		Вісмут ($\ce{Bi}$) & $-165$
	\end{tblr}
\end{table}

\begin{Attention}
	\emphx{Вісмут} --- найсильніший діамагнетик серед звичайних (не надпровідних) речовин: його сприйнятливість на два порядки перевищує
	типову. Причина --- особливості зонної структури (малий ефективний радіус орбіт валентних електронів та специфіка їхнього спектру), що не
	описуються найпростішою атомарною моделлю~\eqref{eq:LangevinDia}, проте не змінюють головного висновку: механізм --- той самий, ларморова
	прецесія електронних орбіт.
\end{Attention}

\subsection*{Ідеальний діамагнетизм: надпровідники}

Окремо варто згадати надзвичайний, \emphx{ідеальний} випадок діамагнітної поведінки --- \emphx{надпровідники}. У надпровідному стані (за
температур, нижчих за критичну) речовина \emphx{повністю виштовхує} магнітне поле зі свого об'єму --- явище, відоме як \emphx{ефект Мейснера}.
Це відповідає значенню $\chi = -1/4\pi$ (тобто $B=0$ всередині зразка за будь-якого зовнішнього поля, аж до критичного значення, вище якого
надпровідність руйнується). На відміну від \enquote{звичайного} діамагнетизму атомних орбіт~\eqref{eq:LangevinDia}, ефект Мейснера --- явище
принципово квантове (пов'язане з незатухаючими надпровідними струмами, що екранують поле), і тому виходить далеко за межі класичної моделі
Ланжевена; тут ми лише фіксуємо цей факт як яскраву ілюстрацію того, наскільки сильним --- за належних умов --- може бути діамагнітний відгук.

%% --------------------------------------------------------
\section{Природа парамагнетизму}
%% --------------------------------------------------------

За аналогією з орієнтаційною поляризацією полярних діелектриків (гл.~\ref{FileldInDielectrics}, п.~\ref{sec:Langevin-Debye}), розглянемо тепер
речовину, атоми якої мають власний ненульовий магнітний момент $p_0$.

За відсутності поля тепловий рух орієнтує атомні моменти цілком хаотично, і сумарна намагніченість дорівнює нулю. У зовнішньому полі $\Bfield$
кожен момент прагне розвернутись уздовж поля (мінімізуючи потенціальну енергію $U=-\vect p_0\cdot\Bfield$, формула~\eqref{eq:PotentialEnergyDipole}),
але тепловий рух цьому протидіє. Рівноважний розподіл орієнтацій установлюється за розподілом Больцмана: імовірність напрямку з енергією
$U=-p_0B\cos\theta$ пропорційна $\exp(p_0B\cos\theta/kT)$. У практично важливому випадку слабких полів і не надто низьких температур
($p_0B\ll kT$) --- а це виконується майже завжди, аж до гелієвих температур і сильних лабораторних полів, --- усереднення дає малий, лінійний
за полем ефект, цілком аналогічний формулі Ланжевена--Дебая~\eqref{eq:LangevinDebye} для полярних діелектриків:
\begin{equation}\label{eq:ParaAvgMoment}
	\langle p_z\rangle = \frac{p_0^2 B}{3kT}.
\end{equation}
Звідси намагніченість $n$ атомів в одиниці об'єму дорівнює $M = n\langle p_z\rangle$, а парамагнітна сприйнятливість --- обернено пропорційна
температурі:
\begin{equation}\label{eq:CurieLawMag}
	\tcbhighmath{
		\chi_{\text{пара}} = \frac{np_0^2}{3kT} \equiv \frac{C}{T},
	}
\end{equation}
де сталу $C=np_0^2/3k$ називають \emphx{сталою Кюрі}, а саму формулу --- \emphx{законом Кюрі} для парамагнетиків.

\begin{Attention}
	На відміну від діамагнітної сприйнятливості~\eqref{eq:LangevinDia}, парамагнітна $\chi_{\text{пара}}$ \emphx{спадає з температурою} (сильніший
	тепловий рух руйнує впорядкування) і, як правило, на \emphx{один--два порядки перевищує} діамагнітний внесок того самого атома при кімнатній
	температурі ($\chi_{\text{пара}}\sim 10^{-3}$--$10^{-4}$ проти $\chi_{\text{діа}}\sim 10^{-5}$--$10^{-6}$). Тому в парамагнітних речовинах
	діамагнітний внесок~\eqref{eq:LangevinDia}, хоча він і присутній завжди, практично непомітний на тлі значно сильнішого парамагнітного відгуку
	--- і лише у речовин без власного атомного моменту (тобто саме в діамагнетиках) він виходить на перший план.
\end{Attention}

Феромагнетики --- це, по суті, парамагнетики з дуже сильною додатковою (обмінною) взаємодією між сусідніми моментами; вище певної критичної
температури (\emphx{точки Кюрі}) ця взаємодія \enquote{програє} тепловому руху, впорядкування руйнується, і феромагнетик перетворюється на звичайний
парамагнетик, що також підпорядковується закону~\eqref{eq:CurieLawMag} (точніше, його узагальненню --- закону Кюрі--Вейса). Детальний розгляд
цього переходу --- предмет окремого розділу.

%% --------------------------------------------------------
\section{Вектор намагніченості}
%% --------------------------------------------------------

Щоб описати стан намагніченого магнетика кількісно, вводять, за прямою аналогією з вектором поляризації~\eqref{eq:Pdef}, \emphx{вектор
намагніченості}, що дорівнює магнітному моменту одиниці об'єму речовини:
\begin{equation}\label{eq:MagnetizationDef}
	\tcbhighmath{
		\Mfield = \frac1V\sum_i \vect p_{m,i},
	}
\end{equation}
де сумування ведеться за магнітними моментами всіх атомів чи молекул $\vect p_{m,i}$, що потрапляють у фізично малий об'єм $V$.

\subsection*{Молекулярні струми як джерело поля намагніченого зразка}

З'ясуємо, як намагнічений зразок створює поле на макроскопічному рівні. Уявімо циліндр, однорідно намагнічений уздовж своєї осі
($\Mfield=\const$), і подумки \enquote{зберемо} його з великої кількості однакових атомних струмових контурів (рис.~\ref{pic:SolenoidEquivalence}),
щільно вкладених один в один, так що момент кожного узгоджений з напрямком $\Mfield$.

\begin{figure}[h!]\centering
	\begin{tikzpicture}[>=latex, scale=1]
		\foreach \x in {0,...,5} {
			\foreach \y in {0,...,2} {
				\draw[thick, blue!60!black] ({1.2*\x},{0.9*\y}) circle (0.42);
				\draw[->, thick, blue!60!black] ({1.2*\x-0.42*0.707},{0.9*\y+0.42*0.707}) arc[start angle=135, end angle=90, radius=0.42];
			}
		}
		\draw[very thick, red] (-0.55,-0.55) rectangle (6.35,2.35);
		\draw[->, very thick, red] (-0.55,-0.55) -- (2.9,-0.55);
		\node[red, right] at (2.9,-0.9) {\small еквівалентний поверхневий струм $i'$};
		\draw[->, thick, brown] (6.9,-0.2) -- (6.9,2.1);
		\node[brown, above] at (6.9,2.15) {\small $\Mfield$};
	\end{tikzpicture}
	\caption{Однорідно намагнічений зразок як набір щільно вкладених атомних струмових контурів: у товщі
		зразка сусідні струми взаємно компенсуються, а некомпенсованим лишається лише
		\emphx{поверхневий} струм $i'$, --- зразок еквівалентний соленоїду\label{pic:SolenoidEquivalence}}
\end{figure}

У товщі зразка струм кожного контуру на межі із сусіднім контуром напрямлений \emphx{назустріч} струму сусіда --- і ці внутрішні струми взаємно
компенсуються, точнісінько так само, як компенсувались об'ємні заряди сусідніх диполів у поляризованому діелектрику
(рис.~\ref{pic:PolarizedParallelepiped}, гл.~\ref{FileldInDielectrics}). Некомпенсованим лишається лише струм на \emphx{бічній поверхні}
циліндра, --- увесь намагнічений зразок поводиться точнісінько так само, як соленоїд: еквівалентний поверхневий струм, що припадає на одиницю
довжини зразка, дорівнює (у гауссовій системі одиниць, порівняйте з полем соленоїда~\eqref{eq:SolenoidField}, де стоїть $nI$):
\begin{equation}\label{eq:SurfaceBoundCurrentUniform}
	i' = cM.
\end{equation}

\begin{Attention}
	Формула~\eqref{eq:SurfaceBoundCurrentUniform} --- це магнітний аналог формули $\sigma'=P_n$~\eqref{eq:SigmaPrime} для зв'язаного заряду.
	Різниця в тому, що там \emphx{нормальна} складова $\vect P$ давала \emphx{скалярний} поверхневий заряд, а тут \emphx{дотична} (тангенціальна)
	складова $\Mfield$ дає \emphx{поверхневий струм}, спрямований по контуру перерізу зразка (а не вздовж осі), --- це прямий наслідок різної
	геометричної природи диполя-заряду і диполя-струму.
\end{Attention}

\subsection*{Загальний випадок. Циркуляція намагніченості}

Для довільної (не обов'язково однорідної) намагніченості узагальнення проводять точно так само, як ми узагальнювали формулу для $\sigma'$ до
теореми Гаусса для $\vect P$~\eqref{eq:GaussForP-diff}: замість потоку через замкнену поверхню тут природно розглянути \emphx{циркуляцію} вектора
$\Mfield$ вздовж довільного замкненого контуру $L$ --- адже саме циркуляція $\Bfield$, а не потік, пов'язана зі струмом (теорема про циркуляцію,
гл.~\ref{MagFieldVacuum}). Не заглиблюючись у громіздкі викладки, наведемо результат: циркуляція $\Mfield$ уздовж контуру $L$ дорівнює
(з коефіцієнтом $1/c$, за аналогією з теоремою про циркуляцію для $\Bfield$) повному молекулярному струму $I'$, що пронизує контур:
\begin{equation}\label{eq:CirculationM}
	\tcbhighmath{
		\oint\limits_L \Mfield\cdot d\vect\ell = \frac1c I'_{\text{всередині}}.
	}
\end{equation}
Стягуючи контур $L$ у точку (як ми робили це в гл.~\ref{MagFieldVacuum}, переходячи від інтегральної форми теореми про циркуляцію до
диференціальної), отримуємо \emphx{диференціальну форму}:
\begin{equation}\label{eq:CurlMBound}
	\tcbhighmath{
		\rot\Mfield = \frac1c\, \vect j',
	}
\end{equation}
де $\vect j'$ --- об'ємна густина молекулярного (зв'язаного) струму. У частковому випадку однорідної намагніченості $\Mfield=\const$ маємо
$\rot\Mfield\equiv0$, тобто $\vect j'=0$ в товщі зразка, --- увесь молекулярний струм зосереджений на поверхні, що узгоджується з
результатом~\eqref{eq:SurfaceBoundCurrentUniform}.

%% --------------------------------------------------------
\section{Вектор напруженості магнітного поля}
%% --------------------------------------------------------

Магнітне поле в речовині створюється \emphx{обома} типами струмів --- і вільними $\vect j$, і молекулярними $\vect j'$. Тому в теоремі про
циркуляцію $\Bfield$ (гл.~\ref{MagFieldVacuum}, універсальний закон, справедливий для будь-яких струмів) слід урахувати обидва внески:
\begin{equation*}
	\rot\Bfield = \frac{4\pi}c\left(\vect j+\vect j'\right).
\end{equation*}
Скориставшись щойно знайденим співвідношенням~\eqref{eq:CurlMBound}, а саме $\vect j'=c\,\rot\Mfield$, маємо
\begin{equation*}
	\rot\Bfield = \frac{4\pi}c\vect j + 4\pi\,\rot\Mfield
	\quad\Longrightarrow\quad
	\rot\left(\Bfield - 4\pi\Mfield\right) = \frac{4\pi}c\vect j.
\end{equation*}
Вираз у дужках позначають окремим символом і називають \emphx{вектором напруженості магнітного поля}:
\begin{equation}\label{eq:HDef}
	\tcbhighmath{
		\Hfield = \Bfield - 4\pi\Mfield.
	}
\end{equation}
Тоді теорема про циркуляцію \emphx{в речовині} набуває вигляду
\begin{equation}\label{eq:CurlHFree}
	\tcbhighmath{
		\rot\Hfield = \frac{4\pi}c\vect j, \qquad \oint\limits_L \Hfield\cdot d\vect\ell = \frac{4\pi}c\sum_i I_i,
	}
\end{equation}
де в правій частині --- \emphx{лише вільні} струми.

\begin{Attention}
	До циркуляції вектора $\Hfield$ входять \emphx{тільки вільні} струми $I_i$ --- унесок молекулярних струмів уже враховано, він
	\enquote{захований} усередині означення~\eqref{eq:HDef}. Це робить $\Hfield$ надзвичайно зручним для розрахунків: щоб знайти циркуляцію
	$\Hfield$, достатньо знати лише вільні струми (наприклад, струм в обмотці котушки), \emphx{незалежно} від того, який магнетик всередині
	котушки.
\end{Attention}

\begin{Attention}
	Зверніть увагу на \emphx{знак}: $\Dfield=\Efield+4\pi\vect P$ (додавання, формула~\eqref{eq:Ddef}), але $\Hfield=\Bfield-4\pi\Mfield$
	(віднімання). Це не помилка, а прояв глибшої відмінності природи двох явищ. Вектор $\Dfield$ будують так, щоб \emphx{включити} внесок
	зв'язаних зарядів у джерело поля через дивергенцію (обидва внески --- $\Efield$ від вільних і від зв'язаних зарядів --- складаються за
	теоремою Гаусса). Натомість поле $\Bfield$ вже \emphx{автоматично} містить внесок від \emphx{усіх} струмів, включно з молекулярними (бо
	теорема про циркуляцію для $\Bfield$ універсальна); тому, щоб отримати величину, циркуляція якої залежить \emphx{тільки} від вільних струмів,
	внесок молекулярних струмів (а він, як ми з'ясували, дорівнює $4\pi\Mfield$ у роторі) доводиться саме \emphx{відняти}. Це та сама
	термінологічна асиметрія \enquote{індукція--напруженість}, про яку йшлося у Увазі після формули~\eqref{eq:BDef} (гл.~\ref{MagFieldVacuum}): у
	вакуумі, де $\Mfield=0$, різниця між $\Bfield$ і $\Hfield$ зникає.
\end{Attention}

\subsection*{Магнітна сприйнятливість і проникність}

Для \emphx{лінійних} магнетиків (переважна більшість дія- і парамагнетиків у не надто сильних полях) намагніченість пропорційна полю
$\Hfield$ \emphx{у самій речовині}, --- за прямою аналогією з $\vect P=\chi\Efield$~\eqref{eq:PlinearE}:
\begin{equation}\label{eq:MlinearH}
	\tcbhighmath{
		\Mfield = \chi\Hfield,
	}
\end{equation}
де $\chi$ --- \emphx{магнітна сприйнятливість} речовини: $\chi<0$ для діамагнетиків~\eqref{eq:LangevinDia}, $\chi>0$ для
парамагнетиків~\eqref{eq:CurieLawMag}. Підставляючи~\eqref{eq:MlinearH} у означення~\eqref{eq:HDef}, отримуємо
\begin{equation*}
	\Bfield = \Hfield+4\pi\Mfield = (1+4\pi\chi)\Hfield.
\end{equation*}
Позначаючи
\begin{equation}\label{eq:MuDef}
	\tcbhighmath{
		\mu = 1+4\pi\chi,
	}
\end{equation}
отримуємо просте співвідношення, аналогічне~\eqref{eq:DepsE}:
\begin{equation}\label{eq:BmuH}
	\tcbhighmath{
		\Bfield = \mu\Hfield.
	}
\end{equation}
Величину $\mu$ називають \emphx{магнітною проникністю середовища}. Для діамагнетиків $\mu$ дещо \emphx{менша} одиниці ($\mu\lesssim1$), для
парамагнетиків --- дещо \emphx{більша} ($\mu\gtrsim1$); в обох випадках відхилення від одиниці --- лише в п'ятому--третьому знаку після коми. Для
феромагнетиків $\mu$ може сягати $10^3$--$10^5$, причому сама залежність $B(H)$ --- нелінійна й неоднозначна (гістерезис), --- питання, яке
залишаємо для окремого розділу.

\begin{table}[h!]\small\centering
	\caption{Порівняння електричного й магнітного полів у речовині\label{tab:ElMagCompareMatter}}
	\begin{tabular}{lcc}
		\toprule
		            & Діелектрик & Магнетик \\
		\midrule
		Дипольний момент одиниці об'єму & $\vect P=\chi\Efield$ & $\Mfield=\chi\Hfield$ \\
		Допоміжний вектор & $\Dfield=\Efield+4\pi\vect P$ & $\Hfield=\Bfield-4\pi\Mfield$ \\
		Теорема з допоміжним вектором & $\divg\Dfield=4\pi\rho$ & $\rot\Hfield=\dfrac{4\pi}c\vect j$ \\
		Матеріальна стала & $\epsilon=1+4\pi\chi$ & $\mu=1+4\pi\chi$ \\
		Зв'язок з допоміжним вектором & $\Dfield=\epsilon\Efield$ & $\Bfield=\mu\Hfield$ \\
		\bottomrule
	\end{tabular}
\end{table}

%% --------------------------------------------------------
\section{Граничні умови на межі магнетиків}
%% --------------------------------------------------------

Як і для $\Dfield$ та $\Efield$ (гл.~\ref{FileldInDielectrics}), знайдемо, як пов'язані поля по обидва боки межі розділу двох магнетиків.

\emphx{Нормальна складова $\Bfield$.} Теорема Гаусса $\divg\Bfield=0$~\eqref{eq:GaussMagDiff} --- універсальний закон, що не залежить від
наявності речовини (реальних магнітних зарядів немає ніде). Застосовуючи її до нескінченно малого циліндра на межі розділу так само, як і у
вакуумі, отримуємо
\begin{equation}\label{eq:BC-Bn}
	\tcbhighmath{
		B_{1n}=B_{2n}:
	}
\end{equation}
нормальна складова $\Bfield$ \emphx{неперервна} на будь-якій межі --- на відміну від $H_n$, яка, узагалі кажучи, зазнає стрибка (бо
$H_n=B_n/\mu$, а $\mu$ по різні боки різна).

\emphx{Тангенціальна складова $\Hfield$.} Застосовуючи теорему про циркуляцію~\eqref{eq:CurlHFree} до вузького прямокутного контуру, що
перетинає межу (аналогічно до виведення умови~\eqref{eq:BC2-E}), знаходимо: якщо на межі немає \emphx{вільного} поверхневого струму, тангенціальна
складова $\Hfield$ неперервна;
\begin{equation}\label{eq:BC-Htau}
	\tcbhighmath{
		H_{1\tau}=H_{2\tau} \qquad (\text{за відсутності вільних струмів на межі}),
	}
\end{equation}
а якщо на межі тече вільний поверхневий струм із лінійною густиною $i$, --- тангенціальна складова $\Hfield$ зазнає стрибка $H_{2\tau}-H_{1\tau} =
\frac{4\pi}c\, i$ (тут $i$ береться в напрямку, перпендикулярному до $\Hfield_\tau$, --- точнісінько так само, як вільний струм у теоремі про
циркуляцію в цілому).

\begin{Attention}
	Порівняйте з діелектриком: там \emphx{нормальна} складова $\Dfield$ неперервна за відсутності вільних \emphx{зарядів}, а \emphx{тангенціальна}
	складова $\Efield$ неперервна завжди. Тут --- дзеркально навпаки: \emphx{нормальна} складова $\Bfield$ неперервна \emphx{завжди}, а
	\emphx{тангенціальна} складова $\Hfield$ неперервна лише за відсутності вільних \emphx{струмів}. Ця перестановка ролей --- ще один прояв тієї
	самої асиметрії \enquote{$\Dfield=\Efield+4\pi\vect P$ проти $\Hfield=\Bfield-4\pi\Mfield$}, про яку йшлося вище.
\end{Attention}

%% --------------------------------------------------------
\section{Підсумки розділу}
%% --------------------------------------------------------

\subsection*{Класифікація й мікроскопічна природа}
\begin{itemize}
	\item Магнетик --- будь-яка речовина: магнітні моменти атомів існують завжди (орбітальний і спіновий рух електронів). Класифікація --- за
	      тим, чи компенсуються ці моменти в атомі без поля: діамагнетики ($\vect p_{\text{атом}}=0$), парамагнетики
	      ($\vect p_{\text{атом}}\neq0$, хаотично орієнтовані без поля), феромагнетики (самочинно впорядковані).
	\item \emphx{Діамагнетизм} --- універсальний, слабкий ($\chi\sim10^{-5}$), не залежний від температури ефект, породжений ларморовою
	      прецесією електронних орбіт у полі: $\omega_L=eB/(2m_ec)$, $\chi_{\text{діа}}=-nZe^2\langle r^2\rangle/(6m_ec^2)$~\eqref{eq:LangevinDia}.
	\item \emphx{Парамагнетизм} --- частковий тепловий розворот власних атомних моментів уздовж поля; сприйнятливість спадає з температурою за
	      законом Кюрі: $\chi_{\text{пара}}=C/T$~\eqref{eq:CurieLawMag}, зазвичай на порядки перевищує діамагнітний внесок.
\end{itemize}

\subsection*{Макроскопічний опис}
\begin{itemize}
	\item Намагніченість: $\displaystyle\Mfield=\frac1V\sum_i \vect p_{m,i}$~\eqref{eq:MagnetizationDef}; еквівалентна молекулярним струмам:
	      об'ємний $\vect j'=c\,\rot\Mfield$~\eqref{eq:CurlMBound}, поверхневий $i'=cM$ для однорідного зразка~\eqref{eq:SurfaceBoundCurrentUniform}.
	\item Вектор напруженості: $\Hfield=\Bfield-4\pi\Mfield$~\eqref{eq:HDef}; циркуляція $\Hfield$ визначається \emphx{лише вільними} струмами:
	      $\rot\Hfield=(4\pi/c)\vect j$~\eqref{eq:CurlHFree}.
	\item Лінійні магнетики: $\Mfield=\chi\Hfield$, $\Bfield=\mu\Hfield$, $\mu=1+4\pi\chi$~\eqref{eq:MuDef}.
	\item Граничні умови: $B_n$ неперервна завжди~\eqref{eq:BC-Bn}; $H_\tau$ неперервна за відсутності вільних поверхневих
	      струмів~\eqref{eq:BC-Htau}.
\end{itemize}

\begin{Attention}
	Головний урок цього розділу: магнітна відповідь речовини --- \emphx{за формою} цілком аналогічна поляризації діелектриків ($\Mfield
	\leftrightarrow \vect P$, $\Hfield \leftrightarrow \Dfield$, $\mu \leftrightarrow \epsilon$), але \emphx{за фізичною суттю} --- явище зовсім
	іншої природи: замість зміщення зв'язаних \emphx{зарядів} тут працює перебудова вже наявних \emphx{мікрострумів} --- або через індуковану
	(ларморову) прецесію орбіт (діамагнетизм, є завжди), або через тепловий розворот уже наявних власних моментів (парамагнетизм, є лише за їх
	наявності). Саме ця відмінність і пояснює, чому знак матеріальної сталої в магнетизмі не такий однозначний, як в електростатиці: якщо
	$\epsilon>1$ практично завжди, то $\mu$ може бути як трохи меншою, так і трохи більшою одиниці --- залежно від того, який з двох механізмів
	переважає.
\end{Attention}

\begin{problem}
Оцінити (за порядком величини) діамагнітну сприйнятливість атома гелію, вважаючи, що обидва його електрони рухаються по орбіті радіуса
$r\approx0{,}5$~\AA, а концентрація атомів відповідає нормальним умовам ($n\approx2{,}7\times10^{19}\ \text{см}^{-3}$). Порівняти отриману оцінку
з табличним значенням для гелію.
\end{problem}

\begin{problem}
Довгий соленоїд з осердям із парамагнетика (сприйнятливість $\chi>0$, $\chi\ll1$) містить $n$ витків на одиницю довжини, якими тече струм $I$.
Знайти індукцію $\Bfield$ всередині соленоїда, а також лінійну густину еквівалентного поверхневого молекулярного струму на бічній поверхні
осердя. У який бік --- за струмом обмотки чи проти нього --- тече цей молекулярний струм?
\end{problem}

\begin{problem}\label{prb:DiaSphere}
Однорідна кулька з лінійного магнетика (магнітна проникність $\mu$) внесена у зовнішнє, далеко від кульки однорідне поле $\Hfield_0$.
Використовуючи аналогію з відповідною задачею для діелектричної кульки в однорідному полі (гл.~\ref{FileldInDielectrics}), записати (без
виведення, за прямою заміною $\epsilon\to\mu$) вираз для намагніченості кульки $\Mfield$ через $\Hfield_0$ і $\mu$.
\end{problem}
